\begin{proof}[Proof of Theorem \ref{thm:expectation}]

The crucial step to apply Equation \ref{eq:mean}
to this data is to work out a simple expression 
for the inverse of the $n\times n$ submatrix $\Sigma_{22}$.
We recall that this matrix is symmetrical with all diagonal entries equal to $\sigma_Q^2 + \sigma_\eta^2$ and all off diagonals equal to $\sigma_Q^2$.

We call upon a lemma due to \cite{miller1981inverse}.
Which states that when $A$ is invertible and $B$ is a rank-1 matrix, the inverse of their sum takes the following form:

\begin{equation}
(A+B)^{-1} = A^{-1} - \frac{1}{1 + \text{Trace}(B A^{-1})} A^{-1} B A^{-1} 
\end{equation}

We can decompose $\Sigma_{22}$
into such an $A$ and $B$
by defining $B$ to be the matrix
that takes value $\sigma_Q^2$ everywhere
and $A$ to be a diagonal matrix that takes values
$\sigma_\eta^2$ (along the main diagonal). 
Thus $\Sigma_{22}^{-1}= (A+B)^{-1}$ and we can proceed by applying the lemma.


First we note that $A^{-1}$ is a diagonal matrix 
with all entries on the main diagonal set to 
$\frac{1}{\sigma_\eta^2}$.
Then we note that $B A^{-1}$ 
is an $n \times n$ matrix with all entries 
set to $\frac{\sigma_Q^2}{\sigma_\eta^2}$.
Thus, $\text{Trace}(B A^{-1}) = \frac{n \sigma_Q^2}{\sigma_\eta^2}$.

The matrix $A^{-1} B A^{-1}$ has all entries equal to $\frac{\sigma_Q^2}{\sigma_\eta^4}$,
and thus our desired inverse $\Sigma_{22}^{-1}$ 
can be expressed as follows:
%
$$\Sigma_{22}^{-1} = (A+B)^{-1} 
=
\begin{bmatrix}
    \frac{1}{\sigma_\eta^2}  &    &0 \\
     & \ddots &   \\
    0  & & \frac{1}{\sigma_\eta^2}\\
\end{bmatrix}
- 
\frac{1}{
1 + \frac{n \sigma_Q^2}{ \sigma_\eta^2}}
\begin{bmatrix}
    \frac{\sigma_Q^2}{\sigma_\eta^4}  &  \dots & \dots \\
    \vdots & \ddots & \vdots  \\
    \vdots  & \dots & \ddots \\
\end{bmatrix}
$$
%
Thus each off-diagonal entry takes values
\begin{align}
-\frac{\frac{\sigma_Q^2}{\sigma_\eta^4}}{
1 + \frac{n \sigma_Q^2}{ \sigma_\eta^2}} &=
%
% multiplying by 1
%
-\frac{\frac{\sigma_Q^2}{\sigma_\eta^4}}{
1 + \frac{n \sigma_Q^2}{ \sigma_\eta^2}} 
\cdot \frac{\frac{\sigma_\eta^2}{\sigma_Q^2}}
{\frac{\sigma_\eta^2}{\sigma_Q^2}}\\
% SIMPLIFYING
& = -\frac{\frac{1}{\sigma_\eta^2}}{
\frac{\sigma_\eta^2}{\sigma_Q^2}+ n}\\
% SIMPLIFYING MORE
& = -\frac{1}{
\frac{\sigma_\eta^4}{\sigma_Q^2}+ n \sigma_\eta^2}
\end{align}
and each on-diagonal entry has an additional $\frac{1}{\sigma_\eta^2}$ term that comes from $A^{-1}$.
\begin{equation}
\frac{1}{\sigma_\eta^2}
- \frac{1}{
\frac{\sigma_\eta^4}{\sigma_Q^2}+ n \sigma_\eta^2}
\end{equation}

Now that we know the precise expression for all entries of $\Sigma_{22}^{-1}$,
we can calculate the vector-matrix product 
$\Sigma_{12} \Sigma_{22}^{-1}$.
Because every entry of $\Sigma_{12}$
takes value $\sigma_Q^2$, and because 
every column of $\Sigma_{22}^{-1}$ has the same $n$ values (just in different order),
the product $\Sigma_{12} \Sigma_{22}^{-1}$
is an $n$ dimensional vector, where
all $n$ values are equal:
\begin{align}
\Sigma_{12} \Sigma_{22}^{-1} & = 
\left[
\frac{\sigma_Q^2}{\sigma_\eta^2}
- \frac{n \sigma_Q^2}{
\frac{\sigma_\eta^4}{\sigma^2} + n \sigma_\eta^2
},
\dots
\right] \nonumber \\
% SIMPLIFY
 & = \left[
    \frac{\frac{\sigma_Q^2}{\sigma_\eta^2}
    \left(
    \frac{\sigma_\eta^4}{\sigma_Q^2} + n \sigma_\eta^2 \right) n\sigma_Q^2
    }
    {\frac{\sigma_\eta^4}{\sigma_Q^2} + n \sigma_\eta^2},
\dots
\right]\nonumber \\
%  SIMPLIFY
 & = \left[
    \frac{\sigma_\eta^2 }
    {\frac{\sigma_\eta^4}{\sigma_Q^2} + n \sigma_\eta^2},
\dots
\right] \nonumber \\
%  SIMPLIFY
 & = \left[
    \frac{1}
    {\frac{\sigma_\eta^2}{\sigma_Q^2} + n },
\dots
\right]  \label{eq:quality-mean}
\end{align}
This expression for $\Sigma_{12}\Sigma_{22}^{-1}$,
together with the definition of the conditional expectation 
of a multivariate Gaussian (Equation \ref{eq:mean}) concludes the proof.

\end{proof}


\begin{proof}[Proof of Theorem \ref{thm:variance}]

We can now produce the expression for 
$\Sigma_{12} \Sigma_{22}^{-1} \Sigma_{21}$.
Because every entry of the $1 \times n$ matrix $\Sigma_{12} \Sigma_{22}^{-1}$
takes value
$$\frac{1}{\frac{\sigma_\eta^2}{\sigma_Q^2} + n }$$
and because every entry in the $n \times 1$ matrix $\Sigma_{21}$
takes value $\sigma_Q^2$, 
\begin{equation}
\Sigma_{12} \Sigma_{22}^{-1} \Sigma_{21} =
\frac{n \sigma_Q^2}{\frac{\sigma_\eta^2}{\sigma_Q^2} + n }
\end{equation}
The expression for the conditional variance
$\text{Var}[Q|Y_1,...,Y_n]$ follows:

\begin{align}
\text{Var}[Q|Y_1,...,Y_n] & = \Sigma_{11} - \Sigma_{12} \Sigma_{22}^{-1} \Sigma_{21} \nonumber \\ 
&=
\sigma_Q^2 - \frac{n \sigma_Q^2}{\frac{\sigma_\eta^2}{\sigma_Q^2} + n }
\nonumber \\
%
& = 
\frac{\sigma_Q^2 \left( \frac{\sigma_\eta^2}{\sigma_Q^2} + n \right) 
- n \sigma_Q^2}{
\frac{\sigma_\eta^2}{\sigma_Q^2} + n
} \nonumber \\
% SIMPLIFY
& = 
\frac{\sigma_\eta^2}{
\frac{\sigma_\eta^2}{\sigma_Q^2} + n
} \nonumber \\
% SIMPLIFY
& = 
\frac{1}{
\frac{1}{\sigma_Q^2} + \frac{n}{\sigma_\eta^2}
} \label{eq:quality-variance}
\end{align}
\end{proof}

